\section{CentralHeatPumpSystem }
\label{centralheatpumpsystem}

\subsection{Overview}
\label{overview-007}

The \emph{CentralHeatPumpSystem} object represents a bank of one or more \emph{ChillerHeaterPerformance:Electric:EIR} modules connected to chilled-water, hot-water, and source plant loops. The wrapper receives the three loop flows, stages available modules, allocates the remaining load and flow to each module, mixes module and bypass outlet streams, and reports the assembled useful loads, source transfer, and electric consumption.

The source connection may serve either side of the refrigerant cycle. It receives condenser heat during cooling-dominant operation and supplies evaporator heat during heating-dominant operation. In simultaneous operation a single module can exchange heat with all three loops; the source receives only the residual after useful heat recovery.

\subsection{Operating Modes and Sequential Dispatch}
\label{operating-modes-and-sequential-dispatch}

Six reported operating states are possible:

\begin{itemize}
\item Mode 0: off.
\item Mode 1: cooling only.
\item Mode 2: heating only.
\item Mode 3: balanced heat recovery with no source transfer.
\item Mode 4: cooling-dominant simultaneous cooling and heating, with residual heat rejected to the source.
\item Mode 5: heating-dominant simultaneous cooling and heating, with residual heat extracted from the source.
\end{itemize}

The wrapper caches the current cooling and heating requests from the two load-side plant connections. A single-mode solver is used when only one request is active. When both are active, the simultaneous solver resolves one refrigerant-cycle result and explicitly routes it among all three connections. Results are independent of whether the cooling or heating plant loop supplies the second request.

Modules are staged sequentially in the order entered. Each module schedule controls its availability. After a module is solved, its useful cooling and heating are subtracted from the remaining loads, and its actual connection flows are subtracted from the remaining wrapper flows. A later module therefore cannot consume load or flow already assigned to an earlier one. Different module capacities, performance curves, schedules, and design flows are retained when a bank is heterogeneous.

The mode of each simultaneous module and the cooling- or heating-dominant state of the wrapper are based on signed source heat transfer, not on a reference capacity ratio. Source heat transfer is positive when heat is rejected to the source and negative when heat is extracted from the source.

\subsection{Thermodynamic and Routing Balance}
\label{thermodynamic-and-routing-balance}

For each active module, compressor power delivered to the refrigerant, evaporator heat, false load, and condenser heat satisfy

\begin{equation}
\dot Q_{cond}=\dot Q_{evap}+\dot Q_{false}+\eta_{motor}P.
\end{equation}

The module result is then routed without altering those thermodynamic quantities. With the source sign convention above, every mode satisfies

\begin{equation}
\dot Q_{heat}+\dot Q_{source}=\dot Q_{cool}+\dot Q_{false}+\eta_{motor}P.
\end{equation}

Cooling-only operation routes evaporator heat to chilled water and all condenser heat to the source. Heating-only operation routes all condenser heat to hot water and extracts evaporator heat from the source. In mode 3, chilled water receives the useful evaporator result and hot water receives the complete condenser result. In mode 4, hot water receives its requested useful heat and the rest of the condenser result goes to the source. In mode 5, chilled water retains its useful cooling while the source supplies the remaining evaporator heat required to produce the hot-water result.

The wrapper reports the sums of module useful cooling, useful heating, signed source heat transfer, and compressor electric power. Ancillary power is added once according to its schedule. Heating-only and heating-dominant operation assign ancillary power to the heating end use; the remaining active modes assign it to cooling. A blank ancillary schedule means continuously available. All runtime values and cached simultaneous results are cleared when the plant does not run the component.

\subsection{Connection Flow and Outlet Mixing}
\label{connection-flow-and-outlet-mixing}

The chilled-water connection uses each module's evaporator design flow, the hot-water connection uses its explicit hot-water design flow, and the source connection allows the larger of its source-extraction and source-rejection design flows. Each heat exchanger uses density and specific heat from its connected plant loop, so a glycol source loop is not evaluated with chilled-water or hot-water properties.

If every module performance object in a wrapper specifies variable flow, the wrapper operates as variable flow. If any module specifies constant flow, the wrapper is constant flow. A legacy bank containing both types is accepted with a warning and every module in that bank uses the constant-flow fallback.

For connection \(k\), let \(\dot m_{k,sys}\) be the plant flow through the wrapper and \(\dot m_{k,i}\) and \(T_{k,i,out}\) be the actual flow and outlet temperature of module \(i\). Sequential allocation ensures

\begin{equation}
\sum_i\dot m_{k,i}\leq\dot m_{k,sys}.
\end{equation}

The bypass flow is

\begin{equation}
\dot m_{k,bypass}=\dot m_{k,sys}-\sum_i\dot m_{k,i},
\end{equation}

and the wrapper outlet temperature is

\begin{equation}
T_{k,sys,out}=\frac{\sum_i\dot m_{k,i}T_{k,i,out}+\dot m_{k,bypass}T_{k,sys,in}}{\dot m_{k,sys}}.
\end{equation}

When the connection has no flow or no active module flow, its outlet temperature remains at the inlet temperature. The same mixing calculation is applied independently to chilled water, hot water, and the source fluid using the final connection-specific module results.

\subsection{Plant Capacity and Design Flow Reporting}
\label{plant-capacity-and-design-flow-reporting}

Plant design-capacity queries are connection specific. The chilled-water connection reports useful evaporator cooling capacity. The hot-water connection reports useful condenser heating capacity, including the compressor power fraction delivered to the refrigerant. The source connection reports a positive capacity-magnitude envelope that covers either cooling heat rejection or heating heat extraction; signed source transfer is used only for runtime reporting.

Maximum and optimum capacities are bank totals. Minimum capacity is the smallest independently stageable module load rather than the sum of every module minimum. This preserves the actual staging behavior of identical and heterogeneous banks.
